%======================================================================
\section{The Limits of the Geometric Framework}\label{sec:limits}
%======================================================================

The Projection Theorem is powerful but not omniscient.
Three fundamental limitations bound the geometric framework, and
intellectual honesty demands that we name them clearly.

\subsection{Causality: the geometry cannot choose the wall}

The geometry tells you \emph{how well} you are projecting.
It cannot tell you \emph{whether you are projecting onto the right wall}.
Omitted variable bias (\S\ref{sec:failures}) is invisible to every
geometric diagnostic: the residuals are still perpendicular, $R^2$ still
decomposes via Pythagoras, the hat matrix is still idempotent.
The Scoreboard reads green for a perfectly projected regression on the
wrong column space.

To choose the right wall, you need \emph{causal reasoning}: directed
acyclic graphs, potential outcomes, natural experiments, instrumental
variables. These are structural tools, not geometric ones.
No amount of staring at shadows will tell you whether the flashlight
is aimed at the right surface.

\subsection{The inner product: whose notion of ``closest''?}

Everything in these notes assumed the Euclidean inner product---all
observations weighted equally, all errors independent.
If observations have \textbf{heteroscedastic} errors
($\text{Var}(\varepsilon_i) = \sigma_i^2$) or are
\textbf{autocorrelated},
the inner product is wrong.
The ``perpendicular'' is not truly perpendicular in the correct metric,
and the ``closest point'' is not truly closest.

\textbf{Generalised Least Squares (GLS)} fixes this by replacing the
Euclidean inner product $\ip{\bm{u}}{\bm{v}}=\bm{u}^\top\bm{v}$ with
$\ip{\bm{u}}{\bm{v}}_\Omega = \bm{u}^\top\bm{\Omega}^{-1}\bm{v}$,
where $\bm{\Omega} = \text{Var}(\bm{\varepsilon})$.
The Projection Theorem still applies, but in a different geometry---a
tilted space where ``perpendicular'' has a new meaning:
\[
  \bhat_{\text{GLS}}
  = (\bX^\top\bm{\Omega}^{-1}\bX)^{-1}\bX^\top\bm{\Omega}^{-1}\by.
\]
The moral: the geometry generalises, but only if you choose the right
inner product.

\subsection{Linearity: the column space is flat}

The column space is \emph{flat}---a hyperplane through the origin.
If the true relationship is curved, the best flat projection may be
a poor representation, like pressing a globe onto a table.
Polynomial features enlarge the column space to include some curvature,
but the bias-variance trade-off (Notebook~3's overfitting trap) limits
how far this can go.
Kernel methods, trees, and neural networks define curved ``column
spaces''---function spaces that lie beyond the reach of linear
projection geometry.

The Projection Theorem is not wrong in these settings.  It is simply
not enough.


%======================================================================
\section{Summary: The Geometric Rosetta Stone}\label{sec:rosetta}
%======================================================================

\subsection*{Closing the circle}

In 1801, Carl Friedrich Gauss took 41 noisy astronomical observations and
dropped a perpendicular from them onto the subspace of Keplerian orbits.
The shadow---the nearest point on that subspace---told him where to
point the telescope.  On New Year's Eve, von~Zach found Ceres within
half a degree of the predicted position.

Every time you call \texttt{lm()} in R or
\texttt{sklearn.LinearRegression().fit()} in Python, you perform
the same act that Gauss performed: you cast a shadow of noisy data
onto a subspace of possible fitted values, you measure the
perpendicular gap that remains, and you trust the right angle.

The mathematics has not changed in two centuries.  Only the notation
is different.

\subsection*{The Rosetta Stone}

Every concept in OLS regression has a geometric face and an algebraic
face.  Table~\ref{tab:rosetta} provides the complete correspondence.
Where the Pythagorean decomposition and $R^2$ require demeaned vectors,
this is noted explicitly.

\begin{table}[ht]
\centering
\caption{The Geometric Rosetta Stone.}\label{tab:rosetta}
\small
\begin{tabular}{@{}lll@{}}
\toprule
\textbf{Statistical concept} & \textbf{Geometric interpretation} & \textbf{Formula / characterisation} \\
\midrule
\multicolumn{3}{@{}l}{\textit{Core projection}} \\[2pt]
Fitted values $\yhat$ & Shadow on the wall & $\bH\by$ \\
Residuals $\be$ & Perpendicular gap & $(\bI-\bH)\by$ \\
Coefficients $\bhat$ & Coordinates of the shadow & $(\bX^\top\bX)^{-1}\bX^\top\by$ \\
Normal equations & Orthogonality condition & $\bX^\top\be=\bm{0}$ \\
Intercept $\hat{\beta}_0$ & Ensures $\bar{e}=0$; centres the & $\bm{1}\in\Col(\bX)$ guarantees \\
 & projection around the data mean & $\be\perp\bm{1}$ \\[4pt]
\midrule
\multicolumn{3}{@{}l}{\textit{Decomposition and fit quality (demeaned vectors $\tilde{\by}=\by-\bar{y}\bm{1}$, $\tilde{\yhat}=\yhat-\bar{y}\bm{1}$)}} \\[2pt]
SST $=$ SSR $+$ SSE & Pythagorean theorem & $\norm{\tilde{\by}}^2=\norm{\tilde{\yhat}}^2+\norm{\be}^2$ \\
$R^2$ & Squared cosine of angle & $\cos^2\theta = \norm{\tilde{\yhat}}^2/\norm{\tilde{\by}}^2$ \\
Degrees of freedom & Dimension of the residual & $\dim\bigl(\Col(\bX)^\perp\bigr) = n-p$ \\
 & subspace & \\
Adjusted $R^2$ & $R^2$ corrected for dimension & $1 - \frac{\norm{\be}^2/(n-p)}{\norm{\tilde{\by}}^2/(n-1)}$ \\[4pt]
\midrule
\multicolumn{3}{@{}l}{\textit{Projection operators and diagnostics}} \\[2pt]
Hat matrix $\bH$ & Projection operator & $\bH^2=\bH,\;\bH^\top=\bH$ \\
Leverage $h_{ii}$ & Arm length in column space & $\diag(\bH)$ \\
Cook's distance $D_i$ & Arm length $\times$ surprise & $\propto h_{ii}\cdot e_i^2$ \\
Gauss--Markov (BLUE) & Closest point in subspace & Projection Theorem \\[4pt]
\midrule
\multicolumn{3}{@{}l}{\textit{Structure and partial effects}} \\[2pt]
FWL theorem & Sequential orthogonal projections & $\bhat_2 = (\bX_2^\top\bM_1\bX_2)^{-1}\bX_2^\top\bM_1\by$ \\
Multicollinearity & Thin, nearly degenerate column space & $\kappa(\bX^\top\bX)\gg 1$ \\
OVB & Wrong column space & $\text{Bias} = \bm{\Gamma}\bm{\beta}_2$ \\[4pt]
\midrule
\multicolumn{3}{@{}l}{\textit{Regularisation (constrained optimisation: trading bias for stability)}} \\[2pt]
Ridge ($L^2$) & Constrained optimisation within a sphere; & Shrinks unstable eigenvector \\
 & trades bias for variance reduction & directions: $\lambda_k/(\lambda_k+\lambda)$ \\
LASSO ($L^1$) & Constrained optimisation within a diamond; & Corner solutions $\Rightarrow$ exact zeros \\
 & trades bias for sparsity & $\Rightarrow$ automatic variable selection \\[4pt]
\midrule
\multicolumn{3}{@{}l}{\textit{Inference}} \\[2pt]
$F$-test & Gap between two projections & $F = \frac{\norm{\yhat-\yhat_R}^2/q}{\norm{\be}^2/(n-p)}$ \\
Confidence ellipsoid & Region of plausible shadows & Axes $\propto 1/\sqrt{\lambda_k}$ \\
\bottomrule
\end{tabular}
\end{table}

\subsection*{The story in sixteen lines}

\begin{center}
\fbox{\parbox{0.88\textwidth}{\centering\itshape
Regression is a projection.\\
Coefficients are coordinates of the shadow.\\
The residual is perpendicular to the column space.\\
The intercept guarantees the residuals have mean zero.\\
$R^2$ is the squared cosine of an angle between demeaned vectors.\\
Variance decomposition is Pythagoras applied to a right triangle.\\
Degrees of freedom count the dimensions the residual can roam.\\
Gauss--Markov says: no linear unbiased estimator gets closer.\\
Controlling for a variable means projecting it away first.\\
Leverage measures how far each point reaches into the column space.\\
Influence is leverage times surprise.\\
Multicollinearity flattens the column space to a sliver.\\
Omitted variables tilt the column space to the wrong wall.\\
The $F$-test measures the gap between two nested projections.\\
Ridge and LASSO trade bias for stability---they are not projections.\\
And geometry alone cannot tell you which wall to aim at.
}}
\end{center}

\bigskip

\noindent
Gauss pointed a pencil at the sky and found an asteroid.  You point a
line of code at a dataset and find a pattern.  In both cases, the
method is the same: drop a perpendicular, read the shadow, respect the
right angle.

\medskip

\noindent
The geometry was always there, hiding inside every formula you
memorised and every $p$-value you reported.  These notes did not
create it.  They just turned on the light.

%======================================================================
% References
%======================================================================

\bigskip
\noindent\textbf{Further reading.}
\begin{itemize}[nosep]
  \item S.\,M.\ Stigler, \emph{The History of Statistics: The Measurement of
        Uncertainty before 1900}, Harvard University Press, 1986.
        \textit{Chapter~4 gives a superb account of Gauss, Ceres, and the
        birth of least squares.}
  \item G.\,W.\ Dunnington, \emph{Gauss: Titan of Science}, Mathematical
        Association of America, 2004 (reprint).
        \textit{The definitive biography; see chapters on the Ceres
        prediction and the development of the method of least squares.}
  \item G.\,A.\,F.\ Seber, \emph{Linear Regression Analysis}, Wiley, 1977.
  \item R.\ Davidson and J.\,G.\ MacKinnon, \emph{Estimation and Inference
        in Econometrics}, Oxford University Press, 1993.
  \item R.\ Christensen, \emph{Plane Answers to Complex Questions: The Theory
        of Linear Models}, 4th ed., Springer, 2011.
  \item S.\ Weisberg, \emph{Applied Linear Regression}, 4th ed., Wiley, 2014.
\end{itemize}

\end{document}
